\answer{ex:05:1}{$\ell\approx17\um$; $\approx100$~días: la distribución por el cerebro la hace la ramificación del cable; la difusión solo emborrona cada sitio de liberación en $\sim20\um$.}
\answer{ex:05:2}{$r=ab/(1+G)$, $S=1/(1+G)$ exactamente; $+1.7\%$ en vez de $+20\%$.}
\answer{ex:05:3}{$T^*=1.21\,\text{s}$ para $G=2$ y $0.19\,\text{s}$ para $G=9$: en la práctica $G\sim2$--$4$, supresión de $3$ a $5$ veces.}
\answer{ex:05:4}{$\mathrm{CV}=1/\sqrt{2\lambda\tau_c}\approx9\cdot10^{-4}$ con frecuencia total $\lambda$; una sola neurona necesita $\tau_c=12.5\,\text{s}$.}
\answer{ex:05:5}{$c\approx0.40$ en ambos casos; $\mathrm{CV}=0.050$ frente a $0.158$, $\sqrt{10}$ veces menos. Las frecuencias de cada neurona siguen siendo proporcionales a $a_i$: solo se regula la suma.}
\answer{ex:05:6}{$N_1=\overline{A\vee B}$, $N_2=\overline{A\vee N_1}$, $N_3=\overline{B\vee N_1}$, $N_4=\overline{N_2\vee N_3}$, XOR $=N_5=\overline{N_4}$; cada conmutación tarda $\tau_c\ln2$; el camino de $4$ compuertas, $2.8\,\text{s}$ por XOR.}
\answer{ex:05:7}{(a)~$\tau_\gamma=6/\ln3\approx5.5\,\text{s}$. (b)~$\bar D<4\,\text{s}$ frente a $D<2.2\,\text{s}$: una espera impredecible vuelve más paciente.}
\answer{ex:05:8}{$k_2=1/\tau_d=10\,\text{s}^{-1}$, $k_1=1/\tau_d-1/\tau_\gamma=9.5\,\text{s}^{-1}$; $\tau_\gamma=1/(k_2-mk_1)$: se duplica con $+2.6\%$ (y con $+5.3\%$ el horizonte es infinito).}
